\chapter{完整训练超参表}
\label{appendix:hyperparameters}

为便于复现，本附录列出本报告所有训练运行的完整超参，按证据标签分组。所有超参在配置文件中均显式给出（路径见每段末），训练代码与种子参考第~3~章实现细节。

\section{Stage 5A 主表（PROJECT\_SUPPORTED, RTX 4070）}

\begin{tabularx}{\textwidth}{lX}
\toprule
项 & 值 \\
\midrule
最大轮数 & 30 \\
批量大小 & 256 \\
优化器 & AdamW (\(\beta_1=0.9, \beta_2=0.999, \varepsilon=10^{-8}\)) \\
学习率 & \(10^{-3}\)（不退火，恒定） \\
权重衰减 & \(10^{-4}\) \\
早停耐心值 & 8（依据验证准确率） \\
学习率调度 & 无（恒定） \\
损失 & 标准交叉熵（\(\epsilon=0\)） \\
增强 & 无 \\
设备 & RTX 4070 (Ada sm\_89), CUDA 12.8, PyTorch 2.9.1+cu128 \\
配置文件 & \path{configs/paper/rml2016a_main_table_full_template.yaml} \\
\bottomrule
\end{tabularx}

\section{扩展训练预算（EXTENDED\_BUDGET\_3090, RTX 3090）}

\begin{tabularx}{\textwidth}{lX}
\toprule
项 & 值 \\
\midrule
最大轮数 & 50 \\
批量大小 & 256 \\
优化器 & AdamW \\
学习率 & 初始 \(10^{-3}\) \\
学习率调度 & 5 轮线性 warmup + cosine annealing 至 0 \\
权重衰减 & \(10^{-4}\) \\
早停耐心值 & 15 \\
损失 & 标准交叉熵 \\
增强 & 无 \\
设备 & RTX 3090 (Ampere sm\_86), CUDA 12.8, PyTorch 2.9.1+cu128 \\
模型范围 & \path{cnn1d}, \path{resnet1d}, \path{iq_param_matched}, \path{fusion_iq_stft} \\
配置文件 & \path{configs/paper/rml2016a_extended_budget_3090.yaml} \\
\bottomrule
\end{tabularx}

\section{低信噪比加权交叉熵（LOW\_SNR\_WEIGHTED\_CE\_3090, RTX 3090）}

\begin{tabularx}{\textwidth}{lX}
\toprule
项 & 值 \\
\midrule
基础设置 & 同上 EXTENDED\_BUDGET\_3090 \\
损失 & 分组 SNR 加权交叉熵；权重 \(w_{\text{low}}=2.0,\ w_{\text{mid}}=1.0,\ w_{\text{high}}=0.7\) \\
SNR 边界 & low\(\le -6\) dB；mid \(\in[-4,6]\) dB；high \(\ge 8\) dB \\
模型范围 & \path{cldnn}, \path{fusion_iq_stft} \\
配置文件 & \path{configs/paper/rml2016a_low_snr_weighted_ce_3090.yaml} \\
\bottomrule
\end{tabularx}

\section{扩展复杂度与 CPU 延迟（CONTROLLED\_LATENCY\_EXTENDED, EPYC CPU）}

\begin{tabularx}{\textwidth}{lX}
\toprule
项 & 值 \\
\midrule
设备 & AMD EPYC 7B12, 强制 \texttt{torch.set\_num\_threads(1)} \\
预热迭代 & 50 \\
计时迭代 & 200 \\
批量大小 & 1 与 256（两次单独测量） \\
PyTorch & 2.9.1+cu128（仅作前向传播，CPU 后端） \\
脚本 & \path{scripts/paper/measure_extended_complexity_latency.py} \\
\bottomrule
\end{tabularx}

\section{提案模型 + 增强 + 标签平滑（FUSION\_CLDNN\_STFT\_AUG\_LS\_3090, RTX 3090）}

\begin{tabularx}{\textwidth}{lX}
\toprule
项 & 值 \\
\midrule
基础设置 & 同上 EXTENDED\_BUDGET\_3090 \\
损失 & 标签平滑交叉熵, \(\epsilon=0.1\) \\
增强（仅训练集） & 相位旋转 \(\theta\sim\mathcal{U}(-\pi,\pi)\), prob 0.5；循环时移 \(k\sim\mathcal{U}(-8,+8)\), prob 0.5 \\
模型 & \path{fusion_cldnn_stft}（CLDNN 风格 I/Q 编码器 + STFT 二维分支 + 融合 MLP 头） \\
模型参数量 & 286,331 \\
配置文件 & \path{configs/paper/rml2016a_fusion_cldnn_stft_aug_ls_3090.yaml} \\
\bottomrule
\end{tabularx}

\section{提案模型单干预消融（FUSION\_CLDNN\_STFT\_ABLATION\_3090, RTX 3090）}

\begin{tabularx}{\textwidth}{lX}
\toprule
项 & 值 \\
\midrule
基础设置 & 同上 FUSION\_CLDNN\_STFT\_AUG\_LS\_3090（3 变体共享） \\
变体 \texttt{arch\_only} & 增强 OFF, \(\epsilon=0\)（纯架构升级） \\
变体 \texttt{arch\_aug}  & 增强 ON,  \(\epsilon=0\) \\
变体 \texttt{arch\_ls}   & 增强 OFF, \(\epsilon=0.1\) \\
模型 & \path{fusion_cldnn_stft}（与上一节相同实现） \\
配置文件 & \path{configs/paper/rml2016a_fusion_cldnn_stft_ablation_3090.yaml} \\
\bottomrule
\end{tabularx}

\section{固定划分与训练种子（所有标签共用）}

\begin{tabularx}{\textwidth}{lX}
\toprule
项 & 值 \\
\midrule
固定划分 & \splitname{} 加密保存为 NPZ 数组 \\
划分文件 & \path{data/splits/rml2016a/stratified_by_mod_snr_seed42.npz} \\
划分摘要 & \path{data/splits/rml2016a/stratified_by_mod_snr_seed42_summary.json} \\
分层维度 & 调制方式 \(\times\) 信噪比联合分层（11 \(\times\) 20 = 220 个组） \\
切分比例 & 70\% / 10\% / 20\% (训练 / 验证 / 测试) \\
切分种子 & 42 \\
训练种子 & 42, 2025, 3407（3 个） \\
\bottomrule
\end{tabularx}
